\documentclass{includes/Thesis}

\usepackage{enumitem}
\usepackage{amssymb}
\usepackage{qrcode}
\urlstyle{same}

% Список первого уровня (itemize)
\setlist[itemize,1]{
    label=\textbullet,
    leftmargin=1.75cm,
    itemsep=0\baselineskip,
    topsep=0\baselineskip
}

% Список второго уровня (itemize)
\setlist[itemize,2]{
    label=$\circ$,
    leftmargin=1.75cm,
    itemsep=0\baselineskip,
    topsep=0\baselineskip
}

% Список третьего уровня (itemize)
\setlist[itemize,3]{
    label=\scalebox{0.6}{$\blacksquare$},
    leftmargin=1.75cm,
    itemsep=0\baselineskip,
    topsep=0\baselineskip
}

\sloppy

\addbibresource{parts/refs.bib}

\DeclareFieldFormat{title}{#1}
\DeclareFieldFormat{journaltitle}{#1}
\DeclareFieldFormat{issuetitle}{#1}
\DeclareFieldFormat{maintitle}{#1}
\DeclareFieldFormat{booktitle}{#1}

\renewcommand*{\mkbibnamefamily}[1]{\textnormal{#1}}
\renewcommand*{\mkbibnamegiven}[1]{\textnormal{#1}}
\renewcommand*{\mkbibnameprefix}[1]{\textnormal{#1}}
\renewcommand*{\mkbibnamesuffix}[1]{\textnormal{#1}}

\DeclareFieldFormat{editortype}{\textnormal{#1}}

\AtBeginBibliography{
    \let\itshape\relax
    \let\slshape\relax
}

\AtBeginDocument{
    \renewcommand{\contentsname}{СОДЕРЖАНИЕ}
}

\begin{thesis}[
        institute = ФЕДЕРАЛЬНОЕ ГОСУДАРСТВЕННОЕ АВТОНОМНОЕ ОБРАЗОВАТЕЛЬНОЕ УЧРЕЖДЕНИЕ ВЫСШЕГО ОБРАЗОВАНИЯ «МОСКОВСКИЙ АВИАЦИОННЫЙ ИНСТИТУТ (НАЦИОНАЛЬНЫЙ ИССЛЕДОВАТЕЛЬСКИЙ УНИВЕРСИТЕТ)»,
        faculty = Компьютерные науки и прикладная математика,
        title = Реализация параллельных алгоритмов интерполяции на основе разреженных сеток с квадратичным базисом,
        year = 2026,
        authorGroup = {М8О-411Б-22},
        educationalProgram = 02.03.02 «Фундаментальная информатика и информационные технологии»,
        authorName = И. Н. Ивкович,
        academicTeacherTitle = {Доктор физико-математических наук, доцент},
        academicTeacherName = А. Ю. Морозов,
        academicSupervisorTitle = {Руководитель образовательной программы},
        academicSupervisorName = С. С. Крылов,
        isAcademic,
        UDC = 519.6,
        keywordsRu={РАЗРЕЖЕННЫЕ СЕТКИ; ИНТЕРПОЛЯЦИЯ; КВАДРАТИЧНЫЙ БАЗИС; ПАРАЛЛЕЛЬНЫЕ АЛГОРИТМЫ; ЧИСЛЕННЫЕ МЕТОДЫ; ДИФФЕРЕНЦИАЛЬНЫЕ УРАВНЕНИЯ; GO; FLUTTER},
        keywordsEn={SPARSE GRIDS; INTERPOLATION; QUADRATIC BASIS; PARALLEL ALGORITHMS; NUMERICAL METHODS; DIFFERENTIAL EQUATIONS; GO; FLUTTER},
    ]
    \setAbstractResource{parts/abstract-ru}{parts/abstract-en}

    \setTerminologyResource{parts/terminology}

    \setIntroResource{parts/intro}

    \addChapter{АРХИТЕКТУРА ПРОГРАММНОГО КОМПЛЕКСА И МАТЕМАТИЧЕСКИЕ ОСНОВЫ МЕТОДА РАЗРЕЖЕННЫХ СЕТОК}{parts/chapter1}
    \addChapter{ПРОЕКТИРОВАНИЕ АРХИТЕКТУРЫ И СТРУКТУР ДАННЫХ ПРОГРАММНОГО КОМПЛЕКСА}{parts/chapter2}
    \addChapter{РЕАЛИЗАЦИЯ ПАРАЛЛЕЛЬНЫХ АЛГОРИТМОВ, АНАЛИЗ РЕЗУЛЬТАТОВ И ПОДГОТОВКА ИЛЛЮСТРАЦИЙ}{parts/chapter3}
    \addChapter{ОБОБЩЕНИЕ РЕЗУЛЬТАТОВ И НАПРАВЛЕНИЯ ДАЛЬНЕЙШЕГО РАЗВИТИЯ}{parts/chapter4}

    \setConclusionResource{parts/conclusion}

    \addAppendix{Репозиторий с исходным кодом}{parts/appendix_a}

\end{thesis}
